\chapter{Étude et développement de nouvelles approches pour la modélisation énergétique de bâtiments}

    \begin{bluebox}
        (26/01) : À date d'écriture du rapport intermédiaire, ce chapitre n'a pas encore été
        traité, les prochaines tâches du projet portant davantage sur l'architecture logicielle
        que les aspects de modélisation.
    
        On peut tout de même supposer que celui-ci rapportera des expériences sur plusieurs
        datasets publics, comme ceux cités en \ref{subsec:reseaux-de-neurones} ou encore ceux
        décrits dans~\cite{subhourlysixdatasets}.
    \end{bluebox}
    + de proche en proche selon l'avancée du modèle
    + méthodes auto-régressives
    + capitalisation sur le modèle existant puisque les NN utilisent beaucoup les
    méta-heuristiques comme la PSO actuellement implémentée
    
    
    \section{Couple modèle-dataset1}
    \section{Couple modèle-dataset2}
    \section{Couple modèle-datasetN}
